The image of the nonempty compact convex polytope \(\mathcal K_D(m_D)\) under the linear map \(q\mapsto (r_D^\dagger)^{\mathsf T}q\) is a nonempty compact convex subset of \(\mathbb R\), hence a closed interval \(I_D\).  For fixed \(p\), strong duality applied to the augmented equality matrix gives
\[
\Gamma_{\rho\prec\sigma}(p;m_D)
=\max_{\eta,\xi}
\{m_D^{\mathsf T}\eta+p\xi:
A_D^{\mathsf T}\eta+(r_D^\dagger)^{\mathsf T}\xi
\le g_{\rho,\sigma,D}^{\mathsf T}\}.
\tag{1}
\]
The feasible set in (1) is a fixed rational polyhedron.  A finite linear program has finitely many relevant bases; on the region where a basis is optimal, its value is affine in \((m_D,p)\), with rational coefficients because \(A_D,r_D^\dagger,g\) are rational.  Hence the value is a finite maximum of affine functions of \(p\), so it is finite, convex, continuous, and piecewise affine on \(I_D\).

The nonpositive sublevel set \(\{p:\Gamma(p;m_D)\le0\}\) is closed and convex.  Within an interval it is empty, a point, or a closed interval, including possibly all of \(I_D\).  Its complement therefore has at most two relative boundary points.  At a boundary, continuity gives value zero; primal and dual attainment in (1) give respectively an equality law and a zero-valued menu-weight certificate.  On an adjacent positive piece the corresponding optimal dual basis has positive value.

The positive common denominator and the checked target locality judgments convert the sign of the numerator gap into the sign of \(Q_\sigma^\star-Q_\rho^\star\), as in the dominance theorem.  If the menu augmented by \((r_D^\dagger,p)\) satisfies its baseline-averaged predicate, the one-district projection theorem applied to that augmented stack lifts the boundary law on the original alphabets while keeping all nonfocal districts fixed.  The two stronger sufficient predicates follow from the established hierarchy.  The final limitations follow because the ratio identity and the lifting equality respectively require locality and exact focal projection.